% ============================================================
%  文档类选项（在入口文件的 \documentclass[...]{tongjithesis} 中设置）
%  field=science    理工科 — 章节编号：1 / 1.1 / 1.1.1
%  本课题为计算机科学与技术（理工），使用 science。
% ============================================================

% ============================================================
%  封面信息
%  任务书 / 开题报告 / 中期报告与论文正文共用本文件；
%  三份行政表格中指导教师为手写签字位，不预填。
% ============================================================
\school{计算机科学与技术学院}
\major{计算机科学与技术}
% 待填写：学号、姓名（AGENTS.md 第 1 节「学生 / 学号 / 班级：待填写」）
\student{待填写}{待填写}
\thesistitle{面向代码智能体的开发任务上下文追踪与代码变更溯源系统设计与实现}{}
% 英文题名：仅首词与专有名词首字母大写（撰写规范）
\thesistitleeng{Design and Implementation of a Task-Context Tracking and Code-Change Provenance System for Code Agents}{}
\thesisadvisor{待填写\quad 待填写}
\thesisdate{\the\year}{\the\month}{\the\day}

% 任务书起讫时间（2026.10.26—2027.05.31，共 32 周，与进度安排表的 32 个周次一致）
% 周数由起讫日期自动换算（向上取整），无需手动填写
\taskbookperiod{2026}{10}{26}{2027}{5}{31}

% ============================================================
%  信息说明页（仅论文正文 main.tex 使用，不影响两份行政表格）
% ============================================================

% 成果类型：thesis — 毕业论文（摘要页标题：摘要）
\infotype{thesis}

% 内容简述（300 字以内）
\infoabstract{%
开发者借助代码智能体完成修改、执行与测试已成日常，但对话、变更、工具调用与测试结果分散在不同界面，缺少以开发任务为中心的组织，使“为什么这样改、依据什么验证、中断后如何接上”难以回答。本课题面向单个开发者，研究任务、对话、变更、工具调用与测试的关联与组织方法，设计并实现支持任务上下文追踪与代码变更溯源的 Web 原型系统，并通过典型编程任务评估其在代码理解、修改原因分析、问题定位与任务衔接方面的作用。研究定位为探索性估计研究，结论以效应量与区间表述。
}

% 毕业论文：正文总字数（成果类型为 thesis 时填写）
\infothesiswords{20000}

% 随附资料（每条用 \item 开头）
\infomaterials{
  \item 原型系统源代码、部署与使用说明（含一键启动脚本与离线样例数据）；
  \item 实验脚本、评估数据与复现说明；
  \item 典型编程任务演示视频与答辩幻灯片。
}
